\chapter{Variables y codificación}
\label{ap:variables}

\section{Codificación de la fuente}

La base de la DGE codifica las variables categóricas con enteros y reserva valores
especiales para la ausencia de dato:

\begin{table}[htbp]
\centering
\caption{Códigos de la fuente relevantes para este trabajo.}
\label{tab:codigos}
\small
\begin{tabular}{l p{10.2cm}}
\toprule
Variable & Codificación \\
\midrule
Comorbilidades y signos & $1=$ sí; $2=$ no; $97=$ no aplica; $98=$ se ignora; $99=$ no especificado \\
\texttt{SEXO} & $1=$ mujer; $2=$ hombre; $99=$ no especificado \\
\texttt{TIPO\_PACIENTE} & $1=$ ambulatorio; $2=$ hospitalizado; $99=$ no especificado \\
\texttt{FECHA\_DEF} & Fecha de defunción; \texttt{9999-99-99} indica paciente vivo \\
\texttt{CLASIFICACION\_FINAL} & $1,2,3=$ caso COVID confirmado; $6=$ sospechoso; $7=$ negativo \\
\texttt{ENTIDAD\_RES} & $1$--$32$ entidades federativas; $97$--$99$ desconocido \\
\bottomrule
\end{tabular}
\end{table}

\section{Variables derivadas}

\begin{table}[htbp]
\centering
\caption{Variables construidas y su papel en el análisis.}
\label{tab:derivadas}
\small
\setlength{\tabcolsep}{4pt}
\begin{tabular}{p{3.5cm} p{3.2cm} p{7.4cm}}
\toprule
Variable & Papel & Definición \\
\midrule
\texttt{defuncion} & Desenlace &
$1$ si \texttt{FECHA\_DEF} $\neq$ \texttt{9999-99-99} \\
\texttt{hospitalizado} & Desenlace (no modelado) &
$1$ si \texttt{TIPO\_PACIENTE} $=2$ \\
\texttt{EDAD} & Predictor &
Valores fuera de $[0,120]$ marcados como faltantes \\
\texttt{sexo\_mujer} & Predictor y atributo sensible &
$1=$ mujer, $0=$ hombre, faltante si $99$ \\
\texttt{<comorbilidad>\_bin} & Predictor &
$1$ si el código es $1$; $0$ si es $2$; faltante en otro caso \\
\texttt{<comorbilidad>\_na} & Indicador de calidad &
$1$ si el código original es $97$, $98$ o $99$ \\
\texttt{entidad\_nombre} & Atributo sensible &
Nombre de la entidad de residencia; eje central de la auditoría \\
\texttt{grupo\_edad} & Atributo sensible &
Cortes 0--17, 18--39, 40--59, 60 o más \\
\bottomrule
\end{tabular}
\end{table}

\section{Predictores admitidos y variables excluidas}

Los trece predictores del modelo son la edad, el indicador de sexo y las once
comorbilidades binarizadas: diabetes, EPOC, asma, inmunosupresión, hipertensión,
otra comorbilidad, enfermedad cardiovascular, obesidad, enfermedad renal crónica,
tabaquismo y neumonía.

Quedan excluidas por posterioridad al punto de predicción, y por tanto por riesgo
de fuga de información: \texttt{UCI}, \texttt{INTUBADO} y \texttt{FECHA\_DEF}
--salvo como origen del desenlace-- así como \texttt{TIPO\_PACIENTE} en su papel de
predictor de mortalidad. La entidad de residencia y el grupo etario se excluyen
por una razón distinta: se reservan como atributos sensibles para la auditoría, no
por fuga.

Los indicadores \texttt{\_na} no se emplean como predictores en el modelo
reportado, pero se conservan en el conjunto de datos porque el dato faltante
resulta informativo y su frecuencia difiere entre entidades por un margen amplio
--de \naEntMin\,\% a \naEntMax\,\%--, lo que constituye un mecanismo plausible de
inequidad. El capítulo~\ref{cap:resultados} lo documenta; explotarlo o corregirlo
queda como trabajo futuro.
